\subsection{Unit Vectors and Orthogonality}

Throughout this subsection, $\mathbb{F}$ is $\mathbb{R}$ or $\mathbb{C}$ and $\mathcal{V}$ is an inner-product space over $\mathbb{F}$.

\begin{definition} \label{def:unit-vector-and-normalization} \hlt{Unit Vector and Normalization}\\
A vector $\mathbf{u}\in\mathcal{V}$ is a \hlt{unit vector} if $\norm{\mathbf{u}}=1$. If $\mathbf{v}\neq\mathbf{0}$, then $\frac{\mathbf{v}}{\norm{\mathbf{v}}}$ is called the \hlt{normalization} of $\mathbf{v}$.
\end{definition}

\begin{proposition} \label{prop:normalization-gives-unit-vector} \hlt{Properties of Normalization}\\
For $\mathbf{v}\in\mathcal{V}\setminus\{\mathbf{0}\}$:
\begin{enumerate}[label=(\roman*)]
\setlength{\itemsep}{0pt}
\item $\frac{\mathbf{v}}{\norm{\mathbf{v}}}$ is a unit vector;
\item $\frac{\mathbf{v}}{\norm{\mathbf{v}}}$ is a positive real scalar multiple of $\mathbf{v}$.
\end{enumerate}
\begin{proof}
\begin{enumerate}[label=(\roman*)]
\setlength{\itemsep}{0pt}
\item[$(i)$] Since $\mathbf{v}\neq\mathbf{0}$, $\norm{\mathbf{v}}>0$ by Proposition \ref{prop:derived-norm-basic-properties}(i)--(ii). The identity $\norm{\frac{\mathbf{v}}{\norm{\mathbf{v}}}}=\frac{1}{\norm{\mathbf{v}}}\norm{\mathbf{v}}=1$ holds by Proposition \ref{prop:derived-norm-basic-properties}(iii).
\item[$(ii)$] $\frac{\mathbf{v}}{\norm{\mathbf{v}}}=\lambda\mathbf{v}$ for $\lambda\defeq\frac{1}{\norm{\mathbf{v}}}\in\mathbb{R}_{>0}$ by Proposition \ref{prop:derived-norm-basic-properties}(i)--(ii).
\end{enumerate}
\hlt{Structural Commentary}: Normalization changes scale but preserves the generated one-dimensional subspace.\\
\hlt{Key Insight}: Divide by the length to make the length equal to one.
\end{proof}
\end{proposition}

\begin{definition} \label{def:orthogonal-vectors} \hlt{Orthogonal Vectors and Finite Lists}\\
Vectors $\mathbf{u},\mathbf{v}\in\mathcal{V}$ are \hlt{orthogonal}, written $\mathbf{u}\perp\mathbf{v}$, if $\innerprod{\mathbf{u},\mathbf{v}}=0$.\\
A finite list $(\mathbf{v}_1,\ldots,\mathbf{v}_m)$ in $\mathcal{V}$ is \hlt{mutually orthogonal} if $(\forall\,j,k\in\{1,\ldots,m\})[j\neq k\Longrightarrow\innerprod{\mathbf{v}_j,\mathbf{v}_k}=0]$.
\end{definition}

\begin{proposition} \label{prop:orthogonality-and-zero} \hlt{Orthogonality and the Zero Vector}\\
For $\mathbf{u},\mathbf{v}\in\mathcal{V}$:
\begin{enumerate}[label=(\roman*)]
\setlength{\itemsep}{0pt}
\item $\mathbf{u}\perp\mathbf{v}\Longleftrightarrow\mathbf{v}\perp\mathbf{u}$.
\item $(\forall\,\mathbf{x}\in\mathcal{V})\,(\mathbf{0}\perp\mathbf{x}\land\mathbf{x}\perp\mathbf{0})$.
\item $\mathbf{v}\perp\mathbf{v}\Longrightarrow\mathbf{v}=\mathbf{0}$.
\end{enumerate}
\begin{proof}
\begin{enumerate}[label=(\roman*)]
\setlength{\itemsep}{0pt}
\item[$(i)$] ($\Rightarrow$) If $\mathbf{u}\perp\mathbf{v}$, then $\innerprod{\mathbf{u},\mathbf{v}}=0$ by Definition \ref{def:orthogonal-vectors}. Conjugate symmetry implies $\innerprod{\mathbf{v},\mathbf{u}}=\overline{\innerprod{\mathbf{u},\mathbf{v}}}=0$, so $\mathbf{v}\perp\mathbf{u}$.
\item[$(i)$] ($\Leftarrow$) If $\mathbf{v}\perp\mathbf{u}$, then $\innerprod{\mathbf{v},\mathbf{u}}=0$ by Definition \ref{def:orthogonal-vectors}. Conjugate symmetry implies $\innerprod{\mathbf{u},\mathbf{v}}=\overline{\innerprod{\mathbf{v},\mathbf{u}}}=0$, so $\mathbf{u}\perp\mathbf{v}$.
\item[$(ii)$] $(\forall\,\mathbf{x}\in\mathcal{V})\,\bigl(\innerprod{\mathbf{0},\mathbf{x}}=0\land\innerprod{\mathbf{x},\mathbf{0}}=0\bigr)$ by Proposition \ref{prop:basic-inner-product-identities}(ii).
\item[$(iii)$] If $\mathbf{v}\perp\mathbf{v}$, then $\innerprod{\mathbf{v},\mathbf{v}}=0$ by Definition \ref{def:orthogonal-vectors}, and $\mathbf{v}=\mathbf{0}$ by definiteness.
\end{enumerate}
\hlt{Structural Commentary}: Conjugate symmetry makes orthogonality symmetric, and definiteness isolates $\mathbf{0}$.\\
\hlt{Key Insight}: Conjugate symmetry preserves the equation ``equals zero.''
\end{proof}
\end{proposition}

\begin{proposition} \label{prop:mutually-orthogonal-nonzero-vectors-independent} \hlt{Mutually Orthogonal Nonzero Vectors Are Independent}\\
Mutually orthogonal nonzero vectors $\mathbf{v}_1,\ldots,\mathbf{v}_m\in\mathcal{V}$ are linearly independent.
\begin{proof}
Suppose $\sum\limits_{j=1}^{m}c_j\mathbf{v}_j=\mathbf{0}$ for $c_1,\ldots,c_m\in\mathbb{F}$, and fix $k\in\{1,\ldots,m\}$.\\
By Proposition \ref{prop:basic-inner-product-identities}(i)--(ii) and mutual orthogonality, $0=\innerprod{\mathbf{0},\mathbf{v}_k}=\innerprod{\sum\limits_{j=1}^{m}c_j\mathbf{v}_j,\mathbf{v}_k}=c_k\innerprod{\mathbf{v}_k,\mathbf{v}_k}=c_k\norm{\mathbf{v}_k}^2$.\\
Since $\mathbf{v}_k\neq\mathbf{0}$, $\norm{\mathbf{v}_k}>0$ by Proposition \ref{prop:derived-norm-basic-properties}(i)--(ii), so $c_k=0$. The index $k$ was arbitrary, hence every coefficient vanishes.\\
\hlt{Structural Commentary}: Orthogonality isolates one coefficient at a time.\\
\hlt{Key Insight}: Pair a dependence relation with one vector from the list; every cross term disappears.
\end{proof}
\end{proposition}
